\documentclass[11pt,a4paper]{article}
\usepackage[utf8]{inputenc}
\usepackage[margin=1in]{geometry}
\usepackage{amsmath}
\usepackage{amssymb}
\usepackage{enumitem}
\usepackage{booktabs}
\usepackage{titlesec}
\usepackage{microtype}

\titleformat{\section}{\large\bfseries\sffamily}{}{0em}{}[\titlerule]
\titleformat{\subsection}{\normalsize\bfseries\sffamily}{}{0em}{}

\setlength{\parindent}{0pt}
\setlength{\parskip}{8pt plus 2pt minus 1pt}

\title{\textbf{CAMS Exam: Long Scenario-Based Select-Two Questions}}
\author{\textbf{Advanced AML Training}}
\date{\today}

\begin{document}
	
	\maketitle
	
	\section*{Instructions}
	Each question presents a detailed scenario followed by five options (A through E). \textbf{Exactly two} of the options are correct. Select both correct answers.
	
	\section{Correspondent Banking \& Nested Accounts}
	
	\subsection{1. The Undisclosed Nested Correspondent}
	
	A European tier-1 bank (GlobalTrust) provides USD correspondent banking services to RegionalBank, a mid-sized commercial bank located in a jurisdiction with weak AML enforcement. During an annual compliance review, GlobalTrust's analytics team discovers that RegionalBank has been routing transactions on behalf of CryptoX, an unregulated virtual asset exchange, through the correspondent account. GlobalTrust never approved CryptoX as a downstream respondent. The transactions appear as large round-sum wire transfers (\$750,000 to \$2 million each) with vague descriptions like "consulting fees" or "technology services." The ultimate originators and beneficiaries are not identifiable from the wire data. GlobalTrust's compliance committee must now decide on next steps.
	
	\begin{enumerate}[label=\Alph*.]
		\item GlobalTrust should immediately terminate all correspondent services to RegionalBank without any prior notification to avoid regulatory penalties.
		\item GlobalTrust should file a Suspicious Activity Report (SAR) with its national FIU detailing the undisclosed nested arrangement.
		\item GlobalTrust should issue a formal Request for Information (RFI) to RegionalBank demanding a complete list of all downstream entities accessing the correspondent account.
		\item GlobalTrust should ignore the finding because RegionalBank is a licensed commercial bank in its home jurisdiction, which is sufficient due diligence.
		\item GlobalTrust should proactively contact CryptoX directly to request their customer transaction data without involving RegionalBank.
	\end{enumerate}
	
	\textbf{Correct Answers: B and C}
	
	\subsection{2. The Payable-Through Account Trap}
	
	A U.S. correspondent bank, First National, enters into a payable-through account (PTA) agreement with a foreign respondent bank, Pacific International Bank (PIB), located in a FATF grey-listed jurisdiction. Under the PTA, PIB's customers can initiate transactions directly through First National's payment systems. Six months later, First National's monitoring detects that multiple PTA accountholders are sending consistent \$9,500 wire transfers to a single shell company registered in the Cayman Islands. The accountholders are all different individuals, but the wires occur within minutes of each other on the same business days. The total aggregate flow exceeds \$5 million over 90 days. None of these accountholders were individually subject to CDD by First National at onboarding.
	
	\begin{enumerate}[label=\Alph*.]
		\item First National has no obligation to investigate because the PTA agreement shifts all AML responsibility to PIB.
		\item First National must apply CDD requirements directly to the underlying PTA accountholders under U.S. regulations.
		\item The pattern of \$9,500 wires suggests possible structuring to evade CTR thresholds and requires immediate investigation.
		\item First National should close the PTA immediately without any notice to PIB or law enforcement.
		\item First National should ignore the pattern because \$9,500 is below the \$10,000 CTR threshold for each individual wire.
	\end{enumerate}
	
	\textbf{Correct Answers: B and C}
	
	\section{Real Estate Money Laundering}
	
	\subsection{3. The Luxury Condo Layering Scheme}
	
	A real estate developer in Manhattan is selling a \$6.5 million luxury condominium unit. The buyer is represented by a law firm that insists on keeping the true buyer's identity confidential, citing attorney-client privilege. The purchase funds arrive as a single wire transfer from a Delaware LLC that was formed only 10 days before the closing. The Delaware LLC's registered agent is a corporate service provider known for forming anonymous companies. The wire originated from a pooled attorney trust account (IOLTA) at a regional bank. The law firm refuses to provide any documentation on the source of funds or the beneficial owner of the LLC. The developer's compliance officer must assess the AML risks.
	
	\begin{enumerate}[label=\Alph*.]
		\item The transaction is exempt from any AML scrutiny because a licensed attorney is involved in the real estate closing.
		\item The combination of anonymous LLC, IOLTA pooled funds, attorney refusal to identify beneficial owner, and rapid LLC formation before closing constitutes multiple red flags requiring enhanced due diligence.
		\item The developer should proceed with the sale because the amount is below the \$10 million threshold for FinCEN's Geographic Targeting Orders.
		\item The developer should consider filing a Suspicious Activity Report with FinCEN if no legitimate explanation is provided after reasonable inquiry.
		\item The developer's only obligation is to verify that the law firm holds a valid professional license.
	\end{enumerate}
	
	\textbf{Correct Answers: B and D}
	
	\subsection{4. The Hotel Purchase Integration Scheme}
	
	A narcotics trafficking network generates \$15 million in drug proceeds over three years. The funds are structured into multiple bank accounts in amounts just below \$9,500 over 24 months. The funds are then aggregated and wired to a newly formed LLC that purchases a beachfront resort property in Florida for \$14 million in cash. The resort operates for two years, generating legitimate hotel revenue. The criminal organization then sells the resort to an arms-length buyer for \$18 million. The sale proceeds are deposited into a commercial bank account and invested in a diversified stock portfolio. A compliance analyst at the bank receiving the \$18 million sale wire flags the transaction for review.
	
	\begin{enumerate}[label=\Alph*.]
		\item The original \$9,500 structuring deposits represent the placement stage of money laundering.
		\item The purchase of the resort with structured funds represents the integration stage of money laundering.
		\item The sale of the resort for \$18 million and deposit of funds represents the integration stage, with the legitimate hotel operations used to "clean" the illicit origin.
		\item The bank receiving the \$18 million sale wire has no obligation to investigate because the sale is between arms-length parties with a legitimate resort operation.
		\item The rapid appreciation from \$14 million to \$18 million in two years, combined with the original anonymous LLC purchase and structured funding, warrants a SAR filing.
	\end{enumerate}
	
	\textbf{Correct Answers: A and C}
	
	\section{Virtual Assets and VASPs}
	
	\subsection{5. The Monero Mixer to Fiat Wash}
	
	A customer of CryptoExchange, a registered VASP, deposits 500 Monero (XMR), an anonymity-enhanced cryptocurrency, over three separate transactions. Within one hour, the customer sends the XMR through a non-regulated decentralized mixer (tumbler) that pools the funds with other users and redistributes them to multiple new wallets controlled by the same customer. After the mixing process, the customer converts the funds to Bitcoin and then immediately requests a fiat wire withdrawal of \$450,000 to a bank account in a jurisdiction known for weak AML controls. When asked about the source of the XMR, the customer provides no documentation and states "it's private personal savings."
	
	\begin{enumerate}[label=\Alph*.]
		\item The use of Monero alone is a standard practice and requires no additional scrutiny.
		\item The sequence of Monero deposit, mixing, conversion to Bitcoin, and rapid fiat withdrawal constitutes a classic layering pattern requiring EDD.
		\item The customer's refusal to explain the source of funds after mixing is a significant red flag warranting potential SAR filing.
		\item The VASP should process the withdrawal immediately to avoid customer dissatisfaction.
		\item The VASP has no obligation to investigate because cryptocurrency transactions are anonymous by design under FATF guidance.
	\end{enumerate}
	
	\textbf{Correct Answers: B and C}
	
	\subsection{6. The NFT Wash Trading Scheme}
	
	A digital artist creates a series of 50 non-fungible tokens (NFTs) on a popular blockchain. Over a 30-day period, two anonymous wallet addresses engage in a pattern of buying and selling the same NFTs back and forth at escalating prices. Wallet A buys an NFT from Wallet B for \$10,000. Two hours later, Wallet B buys a different NFT from Wallet A for \$15,000. This circular trading continues, with the prices artificially rising to \$250,000 per NFT. Blockchain analysis reveals that both wallets are funded from the same centralized exchange account belonging to Individual X. Individual X then sells one NFT to a third-party collector for \$200,000 in cryptocurrency and converts it to fiat currency.
	
	\begin{enumerate}[label=\Alph*.]
		\item This pattern represents wash trading, a form of market manipulation used to artificially establish value and launder money.
		\item The exchange that funded both wallets should consider filing a SAR based on the circular trading pattern.
		\item NFTs are exempt from all AML regulations because they are considered art, not financial instruments.
		\item Individual X's sale to a third-party collector for \$200,000 is a legitimate transaction with no red flags.
		\item The exchange has no obligation to monitor NFT transactions because they occur on a public blockchain outside the exchange's control.
	\end{enumerate}
	
	\textbf{Correct Answers: A and B}
	
	\section{PEPs and Private Banking}
	
	\subsection{7. The Foreign Minister's Adult Child}
	
	A private wealth manager at a Swiss bank is approached by an intermediary representing Ms. Y, the adult daughter of the Minister of Energy of an oil-rich nation with high corruption perception indices. Ms. Y wishes to open an account to deposit \$18 million USD, which she states will be used to purchase luxury real estate in London and Geneva. The intermediary provides documentation showing the funds originate from an offshore company registered in the British Virgin Islands. The company is owned by a trust for which Ms. Y is the sole beneficiary. The stated source of wealth is "consulting fees paid to the offshore company by international energy logistics firms." No further documentation is provided. Ms. Y does not hold any government position herself.
	
	\begin{enumerate}[label=\Alph*.]
		\item Ms. Y should be classified as a PEP family member, requiring Enhanced Due Diligence (EDD) and senior management approval.
		\item The classification as a PEP family member mandates establishing the source of wealth and source of funds through verifiable, corroborative documentation.
		\item Ms. Y is not a PEP because she holds no government position herself, so standard CDD is sufficient.
		\item The ambiguous "consulting fees" from energy logistics firms — directly related to her father's ministry — is a critical red flag requiring intensified scrutiny.
		\item The bank should immediately open the account to capture the \$18 million deposit before the customer goes to a competitor.
	\end{enumerate}
	
	\textbf{Correct Answers: A and D}
	
	\subsection{8. The Former PEP Grace Period}
	
	A private banking client, Mr. A, served as a senior cabinet minister in a foreign government from 2010 to 2018. He retired from public office five years ago and has maintained a banking relationship with your institution for 12 years. During his time in office, the bank applied Enhanced Due Diligence (EDD) and senior management approval. Now, five years after his retirement, a new compliance analyst proposes removing all enhanced monitoring and treating Mr. A as a standard high-net-worth client. The analyst notes that Mr. A's account activity has been consistent and non-suspicious for the entire five-year post-retirement period.
	
	\begin{enumerate}[label=\Alph*.]
		\item Under FATF guidance, PEP status should be maintained for a reasonable period after retirement, typically at least 12 months but potentially longer based on risk.
		\item Five years of clean activity post-retirement is generally sufficient to justify lowering the risk rating, but the decision must be documented and risk-based.
		\item A former PEP can never be reclassified to lower risk under any circumstances.
		\item The bank must consider the country's corruption risk, the former PEP's continued access to state influence, and any ongoing business connections to the government.
		\item The bank is required to maintain EDD on all former PEPs permanently without exception.
	\end{enumerate}
	
	\textbf{Correct Answers: B and D}
	
	\section{Trade Finance and TBML}
	
	\subsection{9. The Grossly Under-Invoiced Electronics Shipment}
	
	A trade finance bank reviews a letter of credit transaction involving the export of 5,000 units of specialized industrial electronics from Country A (high-risk jurisdiction) to Country B (FATF-compliant jurisdiction). The invoice from the exporter, a shell company registered in a secrecy haven, prices each unit at \$80. The global market price for these specific electronic components is \$400 per unit. The shipping documents show the goods were loaded onto a vessel, but no customs inspection was conducted at the origin port. The shipping route includes a transshipment point through a free trade zone known for minimal oversight. The total transaction value is \$400,000 based on the invoice, but the market value would be \$2 million.
	
	\begin{enumerate}[label=\Alph*.]
		\item This transaction exemplifies under-invoicing, where the importer will pay the lower invoice amount but can sell the goods at market price, generating laundered legitimate cash.
		\item The bank should ignore the pricing discrepancy because market prices can fluctuate due to volume discounts and negotiated terms.
		\item The bank should file a Suspicious Activity Report and conduct enhanced due diligence on both the exporter and importer.
		\item The use of a free trade zone transshipment point with minimal oversight is a standard logistics practice with no AML significance.
		\item The bank should process the transaction as normal because the invoice is the only documentary evidence required under international trade rules.
	\end{enumerate}
	
	\textbf{Correct Answers: A and C}
	
	\subsection{10. The Phantom Shipment with Fictitious Invoices}
	
	A manufacturing company in Country X applies for trade finance to purchase "industrial machinery components" from a supplier in Country Y. The supplier is a recently incorporated company with no website, no physical address beyond a mail drop, and no verifiable manufacturing capability. The invoice value is \$2.5 million. The shipping documents include a bill of lading from a carrier that, upon investigation, has no record of the shipment. The bank's trade finance officer notices that the same supplier name appears on three other pending transactions from different buyers, all with identical invoice formatting and similarly vague product descriptions. None of the buyers can provide verifiable end-user certificates.
	
	\begin{enumerate}[label=\Alph*.]
		\item This scenario suggests phantom shipping, where no actual goods exist, and the trade documentation is entirely fabricated to move funds.
		\item The bank should proceed with financing because all documentary requirements (invoice, bill of lading) appear facially valid.
		\item The bank should decline the transaction, file a SAR, and consider terminating the relationship with the buyer.
		\item The bank should contact the supplier directly to request a factory tour before making any determination.
		\item The pattern of the same supplier appearing across multiple suspicious transactions is a critical red flag requiring cross-reference and investigation.
	\end{enumerate}
	
	\textbf{Correct Answers: A and C}
	
	\section{MSBs and Agent Networks}
	
	\subsection{11. The Overnight Volume Spike at a Convenience Store}
	
	National Money Transfer Inc. is a licensed money services business (MSB) with a nationwide agent network of 3,500 retail locations. One of its agents, a small convenience store in a suburban residential neighborhood, has historically processed \$15,000 to \$25,000 in money transfers per month. In the most recent 30-day period, the same agent processed \$1.2 million in outgoing wire transfers — a 5,000\% increase. All transfers are under \$900 (just below the \$900 threshold for certain state reporting requirements), go to the same three foreign countries, and are initiated by customers who appear to be different individuals but all provide the same customer service phone number. The convenience store owner cannot explain the volume increase.
	
	\begin{enumerate}[label=\Alph*.]
		\item National Money Transfer is responsible for its agent network's compliance under BSA regulations and must investigate this anomaly.
		\item The volume spike and pattern of sub-\$900 transfers to the same three countries strongly suggests possible structuring or money mule activity.
		\item National Money Transfer has no liability because each agent is an independent contractor responsible for its own compliance.
		\item National Money Transfer should increase the agent's transaction limits to accommodate the new business volume.
		\item National Money Transfer should consider filing a SAR on the agent location and potentially terminating the agent agreement.
	\end{enumerate}
	
	\textbf{Correct Answers: A and E}
	
	\subsection{12. The Unregistered Check Casher Operating in a Grocery Store}
	
	A state banking regulator receives an anonymous complaint that a grocery store in a major metropolitan area is cashing payroll checks and personal checks for a fee, advertising "Check Cashing — No ID Required — Up to \$5,000." The grocery store is not registered with FinCEN as a Money Services Business. An undercover investigation confirms that the store cashes checks totaling more than \$1,000 for individual customers on a daily basis, maintains a cash inventory of over \$50,000, and does not file Currency Transaction Reports or Suspicious Activity Reports. Some customers are observed cashing multiple checks at the store and immediately purchasing money orders or prepaid cards at the same location.
	
	\begin{enumerate}[label=\Alph*.]
		\item The grocery store is required to register as an MSB with FinCEN because check cashing exceeding \$1,000 per person per day triggers MSB status.
		\item The store's failure to file CTRs for cash transactions over \$10,000 or SARs for suspicious patterns violates the Bank Secrecy Act.
		\item Grocery stores are explicitly exempt from MSB registration requirements under federal law regardless of check cashing volume.
		\item The store's "No ID Required" policy may be permissible for transactions under \$3,000 under certain state laws.
		\item The combination of check cashing, money order sales, and prepaid card sales at the same unregistered location creates a significant money laundering vulnerability.
	\end{enumerate}
	
	\textbf{Correct Answers: A and E}
	
	\section{FATF Effectiveness and Mutual Evaluations}
	
	\subsection{13. The Perfect Laws, Zero Convictions Paradox}
	
	Country Z undergoes a FATF Mutual Evaluation. The assessment team notes that Country Z has enacted comprehensive legislation criminalizing money laundering across all 40 FATF Recommendations. The country has established a Financial Intelligence Unit (FIU) that receives over 50,000 STRs annually from financial institutions. The country has fully implemented the CDD, recordkeeping, and wire transfer rules. However, over the past five years, the country's courts have yielded exactly zero money laundering convictions. The judiciary is known to be severely underfunded, and prosecutors report that judges routinely dismiss ML cases for minor evidentiary technicalities. The country has also never frozen any terrorist assets under UNSCR 1373.
	
	\begin{enumerate}[label=\Alph*.]
		\item Country Z will likely receive a rating of Compliant or Largely Compliant on Technical Compliance because the laws exist on paper.
		\item Country Z will likely receive a rating of Low or Moderate on Effectiveness (specifically Immediate Outcome 7 — ML investigations and prosecutions) due to the zero convictions.
		\item A country cannot receive a low effectiveness rating if it has perfect technical compliance with all Recommendations.
		\item The failure to freeze terrorist assets under UNSCR 1373 will negatively affect the effectiveness rating for Immediate Outcome 10 (Terrorist Financing investigations and prosecutions).
		\item The Mutual Evaluation process only assesses technical compliance, not effectiveness in practice.
	\end{enumerate}
	
	\textbf{Correct Answers: A and D}
	
	\subsection{14. The Grey List Economic Impact}
	
	Jurisdiction X is placed on the FATF "Increased Monitoring" list (commonly known as the Grey List) due to strategic AML/CFT deficiencies. The country has committed to an action plan with FATF but has not yet fully addressed all deficiencies. A major international bank, GlobalCorp, maintains correspondent banking relationships with three banks headquartered in Jurisdiction X. GlobalCorp's risk committee is evaluating whether to maintain, restrict, or terminate these relationships.
	
	\begin{enumerate}[label=\Alph*.]
		\item GlobalCorp is required under FATF standards to terminate all correspondent relationships with Grey List jurisdictions immediately.
		\item GlobalCorp should apply enhanced due diligence (EDD) to all transactions involving Jurisdiction X, including increased monitoring of correspondent accounts.
		\item FATF recommends countermeasures against Grey List jurisdictions, which may include restrictions on business relationships but not automatic termination.
		\item Grey List designation has no practical impact on financial institutions; it is purely a public statement with no required actions.
		\item GlobalCorp should conduct individual, documented risk assessments of each respondent bank rather than applying a blanket termination based solely on the Grey List status.
	\end{enumerate}
	
	\textbf{Correct Answers: C and E}
	
	\section{SAR Filing, Tipping-Off, and Confidentiality}
	
	\subsection{15. The Teller's Inadvertent Disclosure}
	
	A bank's AML analyst, while investigating potential structuring activity, contacts the branch manager and requests copies of customer signature cards and recent deposit tickets for the account of Mr. Smith. The branch teller, overhearing the conversation, notices that Mr. Smith is currently at the teller window attempting to deposit \$9,000 in cash. The teller looks at Mr. Smith and says, "I'm sorry, sir, but our compliance team is investigating your account for structuring patterns. You should probably explain where you get all this cash." Mr. Smith immediately becomes agitated, demands to know what the bank suspects, and then walks out without completing the deposit. No Suspicious Activity Report has been filed yet because the investigation is ongoing.
	
	\begin{enumerate}[label=\Alph*.]
		\item The teller's statement constitutes illegal "tipping-off" because it disclosed the existence of an internal investigation to the customer.
		\item Tipping-off is a violation of AML laws that can result in criminal penalties and regulatory sanctions.
		\item The teller's statement is permissible because no SAR has been filed yet, so no tipping-off has occurred.
		\item The bank should immediately file a SAR documenting the customer's reaction and the teller's disclosure.
		\item The bank should terminate the teller for violating customer service protocols, but no AML violation occurred.
	\end{enumerate}
	
	\textbf{Correct Answers: A and D}
	
	\subsection{16. The Law Enforcement "Keep Open" Request}
	
	A bank has filed three Suspicious Activity Reports over a six-month period on an account held by XYZ Trading Company. The SARs describe a pattern of structured cash deposits just below \$9,500, followed by wire transfers to an offshore shell company in a jurisdiction with weak AML controls. The bank's compliance committee decides to terminate the banking relationship and close the account. Before the account is closed, a federal law enforcement agent contacts the bank in writing and requests that the bank keep the account open for an additional 90 days to allow an ongoing undercover investigation to continue. The agent states that closing the account would alert the subjects and jeopardize the investigation.
	
	\begin{enumerate}[label=\Alph*.]
		\item The bank must close the account immediately regardless of law enforcement's request to minimize its own regulatory risk.
		\item The bank should comply with the law enforcement "keep open" request but must obtain the request in writing and document it thoroughly.
		\item The bank should require senior management approval before agreeing to keep the account open after multiple SAR filings.
		\item The bank should publicly announce that it is cooperating with law enforcement to deter criminal activity on the account.
		\item The bank has no legal obligation to comply with a "keep open" request unless the request is accompanied by a court order or subpoena.
	\end{enumerate}
	
	\textbf{Correct Answers: B and C}
	
	\section{Sanctions and OFAC}
	
	\subsection{17. The Transshipment Evasion Scheme}
	
	A U.S.-based trading company, USA-Import Inc., places an order for industrial equipment from a manufacturer in Germany. The equipment is shipped from Germany to a warehouse in Dubai, where the shipping containers are repackaged with new labels and new bills of lading showing the origin as "UAE." The goods are then shipped to a company in Iran that is designated on the OFAC SDN list. The payments are routed from the Iranian company to a shell company in Turkey, then to the Dubai warehouse operator, and finally to the German manufacturer. USA-Import Inc. pays the German manufacturer via a standard wire transfer from its U.S. bank account. The U.S. bank processing the payment has no knowledge of the Iranian end-user.
	
	\begin{enumerate}[label=\Alph*.]
		\item This typology is known as transshipment sanctions evasion, where the true destination is concealed through intermediate countries.
		\item The U.S. bank processing the payment may be violating OFAC sanctions if it is determined that the bank should have known the ultimate destination was Iran.
		\item The U.S. bank has no liability because the payment was to Germany, a non-sanctioned country, and the bank had no knowledge of Iranian involvement.
		\item USA-Import Inc. has likely violated U.S. sanctions by facilitating the export of goods to Iran without an OFAC license.
		\item The Dubai warehouse operator's actions are legal because transshipment through free trade zones is a standard commercial practice.
	\end{enumerate}
	
	\textbf{Correct Answers: A and D}
	
	\subsection{18. The OFAC 50\% Rule Application}
	
	Company A is owned 60\% by Individual X, who is listed on the OFAC SDN list. Company B is owned 40\% by Company A and 60\% by Individual Y, who is not sanctioned. Company C is owned 30\% by Company A and 30\% by Company B, with the remaining 40\% owned by Individual Z (non-sanctioned). A U.S. bank is processing a wire transfer from Company C to a third-party vendor. The bank's sanctions screening system flags Company C for potential ownership ties to Individual X.
	
	\begin{enumerate}[label=\Alph*.]
		\item Company B is blocked under the 50\% rule because Company A's 40\% ownership of B is aggregated with any other blocked ownership — since Company A is blocked, any entity owned 50\% or more by Company A is blocked, but 40\% is below 50\%, so Company B is NOT automatically blocked.
		\item Company C is blocked if the aggregate ownership by blocked persons (Individual X through Company A and Company B) equals or exceeds 50\%.
		\item The bank should process the wire transfer because Company C is not directly owned by Individual X or any single blocked entity exceeding 50\%.
		\item The bank must trace the ownership through all tiers and calculate the cumulative beneficial ownership by blocked persons to determine the status of Company C.
		\item The 50\% rule applies only to direct ownership and does not require calculation through multiple tiers.
	\end{enumerate}
	
	\textbf{Correct Answers: A and D}
	
	\section{Human Trafficking and Forced Labor}
	
	\subsection{19. The Construction Foreman's Cash Pattern}
	
	A compliance analyst at a regional bank reviews the account of a construction foreman, Mr. Rodriguez. Over the past year, the account has received 150 separate peer-to-peer mobile payments ranging from \$200 to \$600 each from 35 different individuals. All senders have Hispanic surnames. Every Friday, Mr. Rodriguez withdraws 85\% to 95\% of the account balance in cash (typically \$8,000 to \$12,000) and pays utility bills (water, electricity, gas) for three different residential properties with high-density housing. Mr. Rodriguez's stated occupation is "construction foreman," but he files no corporate tax returns for a construction business. The bank has no employment contracts or payroll records linking Mr. Rodriguez to any construction company. The pattern has continued for 18 months.
	
	\begin{enumerate}[label=\Alph*.]
		\item This pattern is highly indicative of human trafficking for forced labor, where multiple victims' wages are deposited to a single controller account.
		\item The payment of utility bills for multiple residential properties suggests that the foreman may be housing trafficking victims and controlling their living expenses.
		\item This is likely a legitimate payroll arrangement for a small construction crew, and the bank should take no action.
		\item The analyst should consider filing a SAR based on the pattern of suspected human trafficking.
		\item The bank should close the account immediately without any investigation because the cash withdrawals appear suspicious.
	\end{enumerate}
	
	\textbf{Correct Answers: A and D}
	
	\subsection{20. The Massage Parlor Real Estate Network}
	
	A regional bank's AML team identifies a network of 12 accounts all linked to the same group of individuals. Each account receives multiple small electronic payments (\$50 to \$150) from dozens of different individuals on a daily basis. The aggregated monthly deposits per account range from \$15,000 to \$30,000. The account holders then transfer most of the funds to a central corporate account owned by a limited liability company that holds title to six commercial properties zoned for "personal services businesses." The LLC then pays mortgages, utilities, and property taxes on these properties. Public records show that the properties operate as massage parlors. The individuals managing the accounts have no verifiable income other than these deposits.
	
	\begin{enumerate}[label=\Alph*.]
		\item This pattern is consistent with human trafficking for commercial sexual exploitation, where illicit proceeds are layered through accounts and used to acquire real estate assets.
		\item The use of an LLC to hold real estate purchased with suspected trafficking proceeds represents the integration stage of money laundering.
		\item The multiple small electronic payments are likely from legitimate customers of a legal massage business.
		\item The bank should file a SAR on the accounts and the LLC based on the trafficking indicators.
		\item The bank has no obligation to investigate because personal services businesses are low-risk customer types.
	\end{enumerate}
	
	\textbf{Correct Answers: A and D}
	
	\section{NPOs and Terrorist Financing}
	
	\subsection{21. The Conflict Zone Charity Diversion}
	
	A non-profit organization registered in Country L (low-risk, FATF-compliant jurisdiction) receives a \$2 million donation from a foreign charity operating in Conflict Zone Z, where terrorist groups are known to operate. The local NPO's stated mission is providing educational supplies and building schools in developing countries. Transaction monitoring shows that within 48 hours of receiving the \$2 million donation, the NPO withdraws \$1.6 million in cash from a branch located 300 kilometers from its registered headquarters, in a region with known terrorist activity. The remaining funds are wired to four other NPOs in neighboring conflict zones. The NPO's leadership cannot provide any documentation of school construction or educational supply purchases matching the cash withdrawals.
	
	\begin{enumerate}[label=\Alph*.]
		\item The rapid cash withdrawal (80\% of funds within 48 hours) in a location distant from the NPO's headquarters is a significant red flag for terrorist financing.
		\item Donations from conflict zones where terrorist groups operate should trigger enhanced due diligence regardless of the donor's stated charitable purpose.
		\item The transaction is legitimate because NPOs often operate in remote areas where cash is the only practical form of payment.
		\item The bank should file a SAR and consider terminating the NPO's account if no legitimate explanation is provided.
		\item NPOs are exempt from AML/CFT monitoring under FATF Recommendation 8 because they are charitable organizations.
	\end{enumerate}
	
	\textbf{Correct Answers: A and D}
	
	\subsection{22. The Shell NPO with No Physical Presence}
	
	A bank's enhanced due diligence team reviews an account for a newly registered non-profit organization called "Global Education Fund." The NPO's website contains generic stock photos, vague mission statements, and no specific projects or past accomplishments. The NPO's incorporation documents list a commercial mail drop as its address. The NPO has no employees on payroll and no physical office. The account receives a \$500,000 wire transfer from a company in Country M, which is on the FATF Grey List. Within one week, the NPO wires \$450,000 to a bank account in Conflict Zone Y, with the memo line stating "humanitarian aid distribution." The remaining funds are withdrawn in cash over several days. The NPO's sole signatory is a student with no prior charitable experience.
	
	\begin{enumerate}[label=\Alph*.]
		\item The NPO's lack of physical presence, employees, office, or verifiable charitable activities indicates it may be a shell entity used for terrorist financing.
		\item The wire to Conflict Zone Y, combined with the origin from a Grey List jurisdiction, warrants immediate SAR filing and potential account termination.
		\item The transaction is legitimate because many NPOs start with minimal infrastructure and operate through volunteers.
		\item The bank should approve the transaction because the memo line states "humanitarian aid distribution," which is a legitimate charitable purpose.
		\item The student signatory's lack of experience is irrelevant to the AML risk assessment.
	\end{enumerate}
	
	\textbf{Correct Answers: A and B}
	
	\section{DeFi and Emerging Technologies}
	
	\subsection{23. The Flash Loan Exploit Layering}
	
	A decentralized finance (DeFi) protocol offers flash loans — uncollateralized loans that must be borrowed and repaid within the same blockchain transaction block. An anonymous user executes a complex flash loan transaction: borrow \$5 million in USDC (a stablecoin), swap the USDC for ETH through a decentralized exchange, deposit the ETH into a liquidity pool, withdraw staked ETH derivatives, swap those back to USDC, and repay the original flash loan — all within three seconds. The net result is \$5 million of USDC that has passed through six different smart contracts, each leaving a record on the blockchain but obfuscating the trail. The user then bridges the funds to a different blockchain and deposits them into a mixing protocol.
	
	\begin{enumerate}[label=\Alph*.]
		\item This flash loan exploit pattern can be used as a layering technique to obfuscate the source of illicit funds through rapid, automated transactions across multiple protocols.
		\item The subsequent transfer of funds to a mixer on a different blockchain further obscures the audit trail and is a high-risk indicator.
		\item Flash loans are inherently legitimate DeFi instruments and cannot be used for money laundering because all transactions are recorded on the blockchain.
		\item DeFi protocols have no AML obligations because they are not centralized financial institutions under any regulatory framework.
		\item A VASP or centralized exchange receiving these funds after the mixer would have limited ability to trace the origin and should apply EDD.
	\end{enumerate}
	
	\textbf{Correct Answers: A and E}
	
	\subsection{24. The Cross-Chain Bridge Privacy Wash}
	
	A cybercriminal converts 100 Bitcoin (BTC), stolen from a regulated exchange, into a privacy token (e.g., RenBTC) and uses a cross-chain bridge to move the funds from the Bitcoin blockchain to a privacy-focused blockchain (e.g., Secret Network). On the privacy blockchain, the funds are deposited into a decentralized mixing contract that splits the 100 BTC into 1,000 separate transactions of 0.1 BTC each, each with different timings and routing through multiple intermediary wallets. After a six-month waiting period, the funds are bridged back to the Ethereum blockchain as Wrapped Bitcoin (WBTC) and deposited into a centralized exchange that has KYC requirements. The user provides identification documents and claims the WBTC originated from "crypto trading profits on a decentralized exchange."
	
	\begin{enumerate}[label=\Alph*.]
		\item The use of a cross-chain bridge to move from a transparent blockchain to a privacy blockchain is a layering technique designed to break the audit trail.
		\item The centralized exchange should treat the deposit with enhanced due diligence, including requests for trading records from the decentralized exchange.
		\item The six-month waiting period before bridging back to a regulated exchange eliminates all money laundering risk.
		\item The exchange should accept the customer's explanation because decentralized exchange trading is legitimate and difficult to document.
		\item The pattern of using privacy tools, mixers, and cross-chain bridges to obscure the origin of funds is a strong indicator requiring a SAR filing if no credible explanation is provided.
	\end{enumerate}
	
	\textbf{Correct Answers: A and E}
	
	\section{Corporate Transparency and UBO}
	
	\subsection{25. The Multi-Tier Bearer Share Structure}
	
	A corporate customer, International Holdings Ltd., applies to open a commercial account at a global bank. The ownership structure is as follows: International Holdings Ltd. is 100\% owned by Caribbean Holdings S.A., a company incorporated in Panama. Caribbean Holdings S.A. is 100\% owned by bearer shares. The bearer share certificate is physically held by a corporate service provider in the Bahamas that serves as a "custodian." The corporate service provider provides a letter stating that the bearer shares are immobilized and that the custodian can identify the bearer share holder upon request. However, the corporate service provider refuses to disclose the identity of the bearer share holder to the bank, citing confidentiality obligations to its client. The beneficial owner of International Holdings Ltd. cannot be identified from any public registry.
	
	\begin{enumerate}[label=\Alph*.]
		\item Under FATF Recommendation 24, countries should prohibit bearer shares or require immobilization with a regulated custodian. However, the bank still needs to identify the beneficial owner for CDD purposes.
		\item The custodian's refusal to disclose the UBO identity means the bank cannot complete adequate CDD and should decline the account opening.
		\item The immobilization of bearer shares with a licensed custodian fully satisfies all beneficial ownership identification requirements.
		\item The bank should accept the custodian's letter as sufficient and open the account without further inquiry into the bearer share holder's identity.
		\item The use of a Panama holding company with immobilized bearer shares is a legitimate corporate structure with no heightened AML risk.
	\end{enumerate}
	
	\textbf{Correct Answers: A and B}
	
\end{document}